\documentclass[12pt]{article}
\usepackage[margin=2.5cm]{geometry}
\usepackage{lmodern}
\usepackage[T1]{fontenc}
\usepackage[italian]{babel}
\usepackage{amsmath,amssymb}
\usepackage{booktabs}
\usepackage{graphicx}
\usepackage{float}
\usepackage{caption}
\usepackage{enumitem}
\usepackage{microtype}

\title{Confronto tra modello a 1 asse (9 regimi) e modello a 2 assi (72 regimi)\\[0.5em]
\large Report tecnico interno}
\author{}
\date{Maggio 2026}

\begin{document}
\maketitle

\section{Obiettivo}

Confrontiamo due specificazioni di un modello AR(1) regime-switching per i prezzi LMP di ISO New England (2021--2025, 7.217 finestre a passo 6\,h):
\begin{itemize}[nosep]
\item \textbf{Modello a 1 asse} (9 regimi economici $E$): $\mu(E)$, $\sigma(E)$, $\varphi(E)$. Matrice di transizione $9 \times 9$.
\item \textbf{Modello a 2 assi} (72 regimi congiunti $E \times D$): $\mu(E)$, $\sigma(E)$, $\varphi(E,D)$. Matrice di transizione $72 \times 72$.
\end{itemize}
L'unica differenza è da dove viene il parametro di mean-reversion $\varphi$: dal solo asse economico (livello di prezzo) o dallo spazio congiunto che include l'asse dinamico (persistenza temporale).

I regimi non sono stimati dal modello: sono quelli prodotti dalla pipeline (ToMATo + Tukey HSD merge) e trattati come dati.

\subsection{Stima dei parametri}

Dato il modello AR(1):
\[
p_t = \mu + \varphi(p_{t-1} - \mu) + \sigma\varepsilon_t, \qquad \varepsilon_t \sim \mathcal{N}(0,1),
\]
i parametri sono stimati come segue:
\begin{itemize}[nosep]
\item $\mu$ e $\sigma$: media e deviazione standard campionaria del prezzo all'interno di ciascun regime $E$. Identici nei due modelli.
\item $\varphi$: coefficiente autoregressivo stimato sulle coppie consecutive di finestre assegnate allo stesso regime:
\[
\hat{\varphi} = \frac{\sum_t (p_t - \mu)(p_{t+1} - \mu)}{\sum_t (p_t - \mu)^2}.
\]
\item \textbf{Modello a 1 asse}: la somma è su tutte le coppie consecutive dentro lo stesso regime $E$. Si ottiene un unico $\varphi$ per regime economico. Esempio: tutte le 1\,024 finestre in $E_5$ (Demand) producono $\hat{\varphi}(E_5) = 0{,}55$.
\item \textbf{Modello a 2 assi}: la somma è su tutte le coppie consecutive dentro la stessa cella $(E, D)$. Si ottiene un $\varphi$ diverso per ogni combinazione. Esempio: le 88 finestre in $E_5 + D_0$ producono $\hat{\varphi} = 0{,}85$; le 230 in $E_5 + D_5$--$D_7$ producono $\hat{\varphi} = 0{,}64$--$0{,}67$.
\item Matrice di transizione $\mathbf{T}$: stimata dalle frequenze relative delle transizioni osservate tra coppie consecutive. $\hat{T}_{ij} = n_{ij} / \sum_k n_{ik}$, dove $n_{ij}$ è il numero di transizioni dallo stato $i$ allo stato $j$. Il modello a 1 asse ha $\mathbf{T}$ di dimensione $9 \times 9$; il modello a 2 assi $72 \times 72$.
\end{itemize}

Tutti i parametri sono stimati sull'intero campione 2021--2025.
Per ciascun modello, generiamo 500 traiettorie simulate della stessa lunghezza dei dati (7\,217 passi).


\section{Risultati aggregati: il paradosso dell'equivalenza}

\subsection{Momenti della distribuzione di $\Delta p$}

La Tabella~\ref{tab:moments} confronta i primi quattro momenti della distribuzione dei cambiamenti di prezzo $\Delta p = p_{t+1} - p_t$ --- empirici e simulati --- insieme alla distanza di Wasserstein.

\begin{table}[H]
\centering
\caption{Momenti di $\Delta p$: distribuzione empirica vs.\ modelli simulati.}
\label{tab:moments}
\begin{tabular}{lrrrrr}
\toprule
 & Media & Std & Asimmetria & Curtosi & Wasserstein \\
\midrule
Empirico       &  0,01 & 28,3 & $-$0,56 & 20,98 & --- \\
$\varphi(E)$   & $-$0,22 & 34,9 &    0,02 &  5,27 & 9,01 \\
$\varphi(E,D)$ & $-$0,43 & 36,3 & $-$0,04 &  5,01 & 10,06 \\
\bottomrule
\end{tabular}
\par\smallskip
\footnotesize\textit{Nota:} Curtosi di Pearson (gaussiana $= 3$). Wasserstein in \$/MWh.
\end{table}

Le differenze tra i due modelli sono trascurabili.
Entrambi sovrastimano la deviazione standard (35--36 vs.\ 28 empirico) e sottostimano pesantemente la curtosi (5 vs.\ 21).
La distanza di Wasserstein è quasi identica (9,0 vs.\ 10,1).

\textbf{Interpretazione:} l'equivalenza aggregata è una proprietà geometrica, non un'indicazione di equivalenza dei modelli.
L'asse $E$ satura la quota dominante della varianza di scala ($\eta^2 = 0{,}82$ sul LMP): i parametri $\mu$ e $\sigma$, identici nei due modelli, determinano la distribuzione marginale.
Il parametro $\varphi$ è adimensionale e controlla la cinematica interna, non la scala.
Su cinque anni, le sovrastime e sottostime locali di $\varphi$ si compensano, producendo l'apparenza di equivalenza.

\subsection{Momenti del livello di prezzo}

\begin{table}[H]
\centering
\caption{Momenti del livello di prezzo $p_t$.}
\label{tab:moments_level}
\begin{tabular}{lrrrr}
\toprule
 & Media & Std & Asimmetria & Curtosi \\
\midrule
Empirico       & 55,6 & 42,1 & 2,36 & 10,12 \\
$\varphi(E)$   & 60,3 & 48,9 & 1,47 &  5,49 \\
$\varphi(E,D)$ & 61,9 & 51,9 & 1,50 &  5,54 \\
\bottomrule
\end{tabular}
\par\smallskip
\footnotesize\textit{Nota:} Curtosi di Pearson (gaussiana $= 3$). Media e Std in \$/MWh.
\end{table}

Anche sui livelli di prezzo i due modelli sono quasi equivalenti.


\section{Coverage: dove emerge la differenza}

\subsection{Coverage globale}

Definiamo la coverage al 90\% come la frazione di prezzi reali che cadono all'interno della banda 5°--95° percentile delle traiettorie simulate a ciascun passo temporale.

\begin{table}[H]
\centering
\caption{Coverage 90\%: globale e condizionata al regime dinamico $D$.}
\label{tab:coverage}
\begin{tabular}{lrrrr}
\toprule
 & Globale & Fast ($D_0$--$D_1$) & Moderato ($D_2$--$D_4$) & Persistente ($D_5$--$D_7$) \\
\midrule
$\varphi(E)$   & 96,3\% & 99,5\% & 97,1\% & \textbf{88,3\%} \\
$\varphi(E,D)$ & 96,8\% & 99,5\% & 97,5\% & \textbf{90,2\%} \\
\bottomrule
\end{tabular}
\end{table}

La coverage globale è quasi identica (96,3\% vs.\ 96,8\%).
La differenza emerge \emph{condizionando} sul regime dinamico:
\begin{itemize}[nosep]
\item Nei regimi \textbf{fast-reverting} ($D_0$--$D_1$): entrambi i modelli coprono il 99,5\%. Nessuna differenza.
\item Nei regimi \textbf{moderati} ($D_2$--$D_4$): differenza minima (97,1\% vs.\ 97,5\%).
\item Nei regimi \textbf{persistenti} ($D_5$--$D_7$): il modello a 1 asse scende a \textbf{88,3\%}, sotto la soglia nominale del 90\%. Il modello a 2 assi raggiunge \textbf{90,2\%}, centrato sul nominale.
\end{itemize}

Questo è il dato chiave.
Il modello a 1 asse usa $\varphi(E_5) = 0{,}55$ per tutte le finestre in $E_5$ (Demand), sottostimando la persistenza ($\varphi$ reale $= 0{,}76$--$0{,}86$) nelle finestre con regime dinamico elevato.
La banda di previsione risulta troppo stretta nei periodi di stress prolungato --- esattamente i periodi in cui una stima accurata è più critica.


\section{Traiettorie simulate}

La Figura~\ref{fig:trajectories} confronta visivamente le traiettorie simulate dai due modelli con il prezzo reale.
Per ciascun modello sono mostrate 5 traiettorie campione (colorate), la mediana simulata (tratteggiata) e la banda al 90\% (sfondo).
Il prezzo reale è in nero.

\begin{figure}[H]
\centering
\includegraphics[width=\textwidth]{results/dual_axis/fig_sim_comparison.png}
\caption{Traiettorie simulate: modello a 1 asse (sopra, rosso) vs.\ 2 assi (sotto, verde). Linea nera: prezzo reale. Banda: intervallo 5°--95° percentile.}
\label{fig:trajectories}
\end{figure}

Le due bande sono visivamente simili --- coerentemente con l'equivalenza aggregata --- ma la banda del modello a 2 assi si adatta leggermente meglio ai periodi di alta volatilità (finestre 1000--1500, corrispondenti all'inverno 2022).


\section{Distribuzione empirica vs.\ simulata}

La Figura~\ref{fig:distributions} sovrappone le distribuzioni empiriche e simulate di $\Delta p$.

\begin{figure}[H]
\centering
\includegraphics[width=\textwidth]{results/dual_axis/report_model_comparison.png}
\caption{Confronto: traiettorie simulate (riga 1) e distribuzioni di $\Delta p$ (riga 2). Sinistra: modello a 1 asse. Destra: modello a 2 assi. Nero: distribuzione/prezzo empirico.}
\label{fig:distributions}
\end{figure}

Come riportato nella Tabella~\ref{tab:moments}, le distribuzioni simulate sono più larghe di quella empirica (Std $\approx 35$ vs.\ 28\,\$/MWh) e mancano le code estreme (curtosi $\approx 5$ vs.\ 21).
Questa è una limitazione dell'AR(1) gaussiano, comune a entrambi i modelli, che non dipende dalla specificazione di $\varphi$.


\section{Sintesi}

\begin{enumerate}
\item \textbf{Sui momenti globali}, i due modelli sono equivalenti. Questo non è un'indicazione di interscambiabilità: è una conseguenza della geometria dello spazio (l'asse $E$ satura la varianza di scala e $\varphi$ è adimensionale).

\item \textbf{Sulla coverage condizionata}, il modello a 2 assi è superiore nei regimi persistenti: 90,2\% vs.\ 88,3\%. La differenza è piccola in valore assoluto ma operativamente significativa: il modello a 1 asse scende sotto la soglia nominale del 90\% proprio nei periodi di stress prolungato.

\item \textbf{Sulle traiettorie}, le bande sono visivamente simili. La differenza è strutturale, non visiva: il modello a 2 assi calibra $\varphi$ sulla condizione dinamica corrente, il modello a 1 asse usa una media.

\item \textbf{Sulla distribuzione di $\Delta p$}, entrambi i modelli sottostimano le code. La limitazione è nell'AR(1) gaussiano, non nella struttura dei regimi.

\item \textbf{Il valore dell'asse $D$} non è nel miglioramento delle metriche aggregate, ma nella \emph{disaggregazione}: rivela che all'interno di ciascun regime economico $E$, il parametro $\varphi$ varia enormemente ($D$ spiega il 64\% della varianza di $\hat{\varphi}$, F-test $p < 10^{-14}$) e che questa variabilità produce una distorsione asimmetrica nella caratterizzazione del rischio.
\end{enumerate}

\end{document}
